

%------------------------------------------------------------------
%	PACKAGES AND OTHER DOCUMENT CONFIGURATIONS
%------------------------------------------------------------------

\documentclass[letterpaper]{inzane_syllabus} % a4paper for A4

\usepackage{booktabs, colortbl, xcolor}
\usepackage{tabularx}
\usepackage{enumitem}
\usepackage{ltablex} 
\usepackage{multirow}

\setlist{nolistsep}

\usepackage{lscape}
\newcolumntype{r}{>{\hsize=0.9\hsize}X}
\newcolumntype{w}{>{\hsize=0.6\hsize}X}
\newcolumntype{m}{>{\hsize=.9\hsize}X}

\renewcommand{\familydefault}{\sfdefault}
\renewcommand{\arraystretch}{2.0}
%-----------------------------
%	 PERSONAL INFORMATION
%-----------------------------

%\profilepic{Course_Graphic.png} % Profile picture, if the height of the picture is less than that of the cirle, it will have a flat bottom. 
%\profilepic{Sarcomere_Structure.png}
\profilepic{Course_Graphic.png}

% To remove any of the following, you need to comment/delete the lines in the .cls file (c. line 186). Commenting/deleting the lines below will produce an error. 

%To add different lines, you will need to create the new command, e.g. \profPhone, in the .cls file c. line 76, and command to create the line in the side bar in the .cls file c. line 186

\classname{Clinical \\ Physiology} 
\classnum{PTH6111.01 (Summer 2025) \\ Dept of Physical Therapy} 

%%%%%%%%%%%%%%% PROF INFO
\profname{Sean Collins, PT, ScD}
\officehours{Office Hrs: MW 1-3pm}
\office{SRH (Hall) 211}
\site{\href{https://mycourses.plymouth.edu/courses/21766}{Canvas}} 
\email{smcollins1@plymouth.edu}

%%%%%%%%%%%%%%% COURSE  INFO
\prereq{Credit Hours: 3}
\classdays{Mon, Weds, Friday}
\classhours{11a-12:15p}
\classloc{HPC W210}

\about{Semester Clock Hours: 40 \\ \hspace{} \\ Weekly Clock Hours: 4.5 \\ \hspace{} \\ Prerequisites: DPT Matriculation \\ Corequisites: PTH6110, PTH6112, PTH6115, PTH6116 \\ \hspace{} \\ Mode of Delivery: In Person \\ \hspace{} \\ Learning Activities: Lecture, Discussion \\ \hspace{} \\Lab Component: 3 hours of class time}

%------------------------------

\begin{document}

%-----------------------------
%	 DESCRIPTION


\makeprofile % Print the sidebar

%------------------------------
%	 OVERVIEW
%------------------------------
\vspace{2 cm}



\section{Course Description}
\vspace{0.5cm}

Clinical Physiology covers the function of the muscular; and structure / function of the neuromuscular, circulatory, pulmonary, endocrine, renal and metabolic systems. Students relate anatomical and physiological systems to functional capabilities. Students apply physiology to reasoning about clinical scenarios and apply concepts to patient/client problems with the WHO International Classification of Function for system organization.

\vspace{0.5cm}
\section{Student Learning Outcomes}
\vspace{0.5cm}
%use \begin{outline} or \begin{outline}[enumerate] to create a list with subitems. 
\begin{enumerate}
\item Describe and discuss the anatomy and function of the cells and cellular components of the muscular, circulatory, pulmonary, endocrine, renal, and metabolic systems. (7A)
\item Describe the development of the musculoskeletal, circulatory, pulmonary, endocrine, renal and metabolic systems. (7A)
\item Describe and analyze the different energy systems (ATP/CP, glycolysis, oxygen transport/ electron transfer) used in different activities (eg, one minute, anaerobic, or aerobic activities). (7A)
\item Analyze systematic interrelationships between the muscular, circulatory, pulmonary, endocrine, renal and metabolic systems in functional activities. (7A)
\item Identify possible causes of perturbations to homeostasis at rest and with physiological challenges (physical stress, nutritional stress, environmental stress). (7A, 7C)
\item Relate system function and capacity with multi system function and capacity. (7A, 7C)
\item Demonstrate the ability to apply the basic concepts of physiology to the analysis of patient/client problems related to the muscular, circulatory, pulmonary, endocrine, renal and metabolic systems systems. (7A, 7C)
\item Relate the the International Classification of Function (ICF) to describe how physiological systems emerge as function and consider the whole system impact of dysfunction (impairments) (7D21)
\item Show how to perform cardiovascular/pulmonary tests and measures including: Heart rate, Respiratory rate, pattern and quality, Blood pressure, Pulse Oximetry in an isolated practice condition. (7D19)
\end{enumerate}

\vspace{0.5cm}

%---------------------------------
%	 COURSE MATERIALS
%---------------------------------
\vspace{0.5cm} 

\section{Course Materials}
\vspace{0.2cm}

{\color{myCOLOR} Required Texts}
\vspace{0.2cm}

Collins SM. Clinical Physiology: A muscle centered approach. Manuscript in Preparation, 2024. 

{\color{myCOLOR} Other}\\
Link to: \href{https://chatgpt.com/g/g-67e58f208c708191b17bc5fb184fae1d-clinical-physiology-course-companion}{Clinical Physiology Course Companion GPT}  \\ 
Additional readings available on Canvas. 

{\color{myCOLOR} References}\\
All references utilized in the book and in the class are available in an online Zotero group library (if you would like access you need to join Zotero (free) and request to be a member of the group) at this link: \href{https://www.zotero.org/groups/4709608/clinical_phys_muscle}{Clinical Physiology Zotero Group}  


%%%%%%%%%%%%%%%%%%%%%%%%%%%%%%%%%%%%%
%                SECOND PAGE
%%%%%%%%%%%%%%%%%%%%%%%%%%%%%%%%%%%%%

\newpage % Start a new page

%\makeSide % Print the FAQ sidebar; To get rid of, simply comment out and uncomment 
\makeFullPage

%----------------------------------------------------------------------------------------
%	 GRADING
%----------------------------------------------------------------------------------------
\vspace{0.5cm}
\section{Grading}
\vspace{0.5cm}

There are a total of 100 points available in this class. These points will come from the following areas.

\begin{twentyshort}
	\twentyitemshort{60 points}{Quizzes (5)}
    \twentyitemshort{35 points}{Cumulative Final Exam}
    \twentyitemshort{5 points}{Vital Sign Competency Check}
    \twentyitemshort{Adjustment}{Professional Conduct (See below)}
\end{twentyshort}

\subsection{Grading Scale}

A $\ge$ 94  \\
90 $\le$ A- < 94 \\
87 $\le$ B+ < 90 \\
83 $\le$ B < 87 \\
80$\le$ B- < 83 \\
77 $\le$ C+ < 80 \\
73 $\le$ C < 76 \\
F < 73 \\

Students are reminded that continuation in the program necessitates achieving an overall grade point average of 3.0 (B) or better.  Grades below 2.0 (C) are not acceptable or counted toward the DPT degree.  No more than 10 credits below 3.0 may be counted toward the DPT degree (grades of B-, C+, or C). 

\section{Assignment Descriptions}
\subsection{Quizzes (60 points total)}

\begin{twentyshort}
	\twentyitemshort{Quiz 1 (10 points)}{Chapters 1 \& 2 (Introduction \& Fundamentals)}
    \twentyitemshort{Quiz 2 (20 points)}{Chapters 3 - 6 (Muscle Function)}
    \twentyitemshort{Quiz 3 (10 points)}{Chapters 7 \& 8 (Microcirculation \& Renal Clearance)}
    \twentyitemshort{Quiz 4 (10 points)}{Chapters 9 \& 10 (Circulation \& Cardiac Pump)}
    \twentyitemshort{Quiz 5 (10 points)}{Chapters 11 \& 12 (Respiration \& Ventilation)}
\end{twentyshort}

Quizzes occur throughout the semester (see the schedule for the exact dates).  The allotted time is 72 seconds per question and the number of questions will be announced prior to the quiz. The remaining class time will be used to review and discuss the quiz. 

\subsection{Cumulative Final Exam (35 points total)}

\begin{itemize}
\item Part 1 based on material from Quiz 1: 5 points towards final exam
\item Part 2 based on material from Quiz 2: 5 points towards final exam
\item Part 3 based on material from Quiz 3: 5 points towards final exam
\item Part 4 based on material from Quiz 4: 10 points towards final exam
\item Part 5 based on material from Quiz 5: 10 points towards final exam
\end{itemize}

\subsection{Quiz Grade Replacement}

Quiz grade replacement, also known as the final exam priority, is that grades on Quiz 1 - 5 can be replaced by a higher score achieved on the respective part of the Final Exam (Parts 1 - 5, respectively). For example, if you earned 6/10 points in Quiz 1, but earned 4/5 points in Part 1 of the Final Exam, your Quiz 1 grade will change to 8/10. This course policy encourages students to continue to improve their understanding throughout the semester and to reflect on that final state of understanding in the final grade. It also negates the need for makeup quizzes throughout the term.
\bigskip
\subsection{Vital Sign Competency Check (5 points total)}

Students will demonstrate competency in the practical skills learned in the two lab sessions of this course (cardiovascular and pulmonary vital signs). There are two lab sessions in the final two class meetings. The first session demonstrates and provides an opportunity to practice the skills. The second session provides an opportunity for each student to demonstrate competency in the skill. If competency is not met in this session, the student must arrange for a competency check before the end of the term.

\subsection{Makeup Quizzes}

There will be no makeup quizzes (see above, including the Quiz Grade Replacement section). A makeup final exam will be given only in specific cases of absence (documented absences). If possible, circumstances that may warrant absence should be discussed with the instructor prior to the scheduled time of the exam. Failure to notify the instructor as stated may result in a zero for the exam.

\subsection{Professional Conduct}

\footnotesize{Students are reminded of the program expectation that they follow Professional Conduct, see the Student Handbook, Professional Conduct and Reasonable Professional Expectations for the process of handling violations, remediation and consequences at the program level. Violations of Professional Conduct (such as unexcused absences) will result in a verbal warning. Continued violation will result in a written notification of a grade reduction of up to 10\%. Violations of Professional Conduct that reach an egregious level, such as cheating on a quiz or exam, will result in immediate failure of that quiz or exam.}
\bigskip

\section{Program \& University Policies \& Information}

\subsection{Doctor of Physical Therapy Student Handbook}
Students are reminded that more detailed information on program policies, including the appropriate and ethical use of Artificial Intelligence (AI), can be found in the Plymouth State University Doctor of Physical Therapy Student Handbook.

\subsection{Academic Standing}
Students are reminded that continuation in the program necessitates achieving an overall grade point average of 3.0 (B) or better.  Grades below 2.0 (C) are not acceptable or counted toward the DPT degree.  No more than 9 credits below 3.0 may be counted toward the DPT degree (grades of B-, C+, or C). Given the Make-up examinations will be given only in specific cases of absence (documented absences) for the Final Exam.  Circumstances that may warrant an absence should be discussed with the instructor before the scheduled time of the exam.  Failure to notify the instructor as stated may result in a zero for the exam.

\subsection{Fair Grading Policy}
\footnotesize{Fair and equitable grading reflects values to which all members of the Plymouth State community commit themselves. Grades are used to assess the relative extent to which students achieve course objectives in all for-credit courses at PSU. Academic freedom allows instructors to determine course objectives within the bounds of established curricula, and the means by which a student’s mastery of those objectives will be evaluated; and evaluate the quality of work on individual exams or assignments. Information can be found at \url{https://coursecatalog.plymouth.edu/university-policies-procedures/} under Fair Grading.}

\subsection{Graduate Academic Policies}
The DPT Program is a selective admissions cohort-structured program.  Some Graduate Academic Policies may not be relevant to the program.  Students should refer to the Doctor of Physical Therapy Program Handbook for specific information about Scheduling, Capstones, and Internships.   The sections on Grading System and Academic Standing Policy do apply to DPT students: \url{https://coursecatalog.plymouth.edu/university-policies-procedures/graduate-academic-policies/}

\subsection{Student Support}
\footnotesize{Student Support Foundation (SSF) provides short-term emergency financial assistance and long-term student support. For more information, see: \url{https://campus.plymouth.edu/student-support-foundation/}  Student Support Foundation also runs a food pantry, located in Belknap Hall. To learn more about SSF or access the food pantry, either via open hours or a private appointment, contact the SSF advisor, at psu-ssf@plymouth.edu }

\subsection{Academic Dishonesty}
\footnotesize{Cheating and plagiarism are very serious forms of academic dishonesty. Any use of unauthorized assistance on exams, papers, homework assignments, or other course work constitutes cheating. Knowingly providing assistance during an exam or allowing other students to copy one’s work is also a serious form of academic dishonesty. Plagiarism consists of submitting written work that has been developed wholly or partially by someone else. Submitting written work in which the ideas of others have been duplicated or even paraphrased without proper reference to the author is also a form of plagiarism. Also considered plagiarism is the acquisition of term papers or other assignments from another source and subsequent presentation of these materials as the student’s own work. In addition, students may not use papers in more than one course without the permission of both instructors. 
 
Any student suspected of academic dishonesty will be reported to the Program Director immediately. A full description of academic integrity, violations and procedures when a violation to academic integrity in a course is suspected can be found at \href{https://coursecatalog.plymouth.edu/university-policies-procedures/}{this website under Academic Integrity}.  
}
\newpage
\subsection{Attendance Policy}
\footnotesize{Attendance is expected at all class sessions. Students are ultimately responsible for their attendance, however faculty will assess students Professional Expectations based, in part, on their attendance (See Professional Conduct under Student Evaluation/Grading).  Students who miss class for any reason are responsible for the information and work missed. Students should speak with the professor before missing the scheduled class when able.}


\subsection{Campus Accessibility Policy}
\footnotesize{Plymouth State University is committed to providing students with documented disabilities equal access to all university programs and facilities. If you think you have a disability that requires accommodations, you should immediately contact Campus Accessibility Services (CAS) in Speare 210 (535-3300) to determine whether you are eligible for such accommodations. Academic accommodations will only be considered for students who have registered with CAS. If you have a Letter of Accommodation for this course from CAS, please provide the instructor with that information privately so that you and the instructor can review those accommodations: \url{https://campus.plymouth.edu/accessibility-services/}}

\subsection{LibGuide for Physical Therapy:}
\footnotesize{The Libguide for Physical Therapy is an information sharing platform created through a collaboration between the library and the DPT faculty that presents relevant resources to support students in the Physical Therapy program. Examples of information you will find include tutorials to better understand citing sources in APA and AMA style, resources on how to avoid plagiarism, quick access to relevant databases and interlibrary loan as well as a contact link to the programs’ assigned librarian. To view the DPT Program Libguide go to \url{https://library.plymouth.edu/physical-therapy}.}

\vspace{0.5cm}
\subsection{Cell Phones}
Cell phone use (including texting) is prohibited in class. Also, please respect your classmates by turning your phone to vibrate prior to class start time. Do not turn phones off as the University uses texts to relate emergency information.  

\subsection{Diversity and Inclusivity Statement}

I consider this classroom to be a place where you will be treated with respect, and I welcome individuals of all ages, backgrounds, beliefs, ethnicities, genders, gender identities, gender expressions, national origins, religious affiliations, sexual orientations, ability - and other visible and non-visible differences. All members of this class are expected to contribute to a respectful, welcoming and inclusive environment for every other member of the class. See this web page for the \href{https://www.plymouth.edu/idea-center}{Plymouth State University IDEA Center
Inclusion, Diversity, Equity, and Access.}

%%%%%%%%%%%%%%%%%%%%%%%%%%%%%%%%%%%%%%%%%%%%%%%%%%%%%%%%%%%%%%%%%%%%%%%%%%%%%
%                COURSE SCHEDULE
%%%%%%%%%%%%%%%%%%%%%%%%%%%%%%%%%%%%%%%%%%%%%%%%%%%%%%%%%%%%%%%%%%%%%%%%%%%%%
%\newpage
%\makeFullPage
\bigskip
\bigskip
\section{Tentative Class Schedule}

Note: As changes arise they will be discussed in class.

There is a quiz debrief following each quiz. 


\begin{table}[h!]
\centering
\begin{tabular}{p{1.9cm}p{4.86cm}p{4.03cm}p{5.71cm}}
\hline
\multicolumn{1}{|p{1.9cm}}{Week of} & 
\multicolumn{1}{|p{4.86cm}}{Monday} & 
\multicolumn{1}{|p{4.03cm}}{Weds} & 
\multicolumn{1}{|p{5.71cm}|}{Friday} \\ 
\hline\hline
\multicolumn{1}{|p{1.9cm}}{June 2} & 
\multicolumn{1}{|p{4.86cm}}{Chp 1 - Introduction} & 
\multicolumn{1}{|p{4.03cm}}{Chp 2 - Fundamentals} & 
\multicolumn{1}{|p{5.71cm}|}{Chp 2 - Fundamentals (cont)} \\ 
\hline
\multicolumn{1}{|p{1.9cm}}{June 9} & 
\multicolumn{1}{|p{4.86cm}}{Quiz 1: Chps 1 \& 2 } & 
\multicolumn{1}{|p{4.03cm}}{Chp 3 - Tension} & 
\multicolumn{1}{|p{5.71cm}|}{Chp 4 - Excitation } \\ 
\hline
\multicolumn{1}{|p{1.9cm}}{June 16} & 
\multicolumn{1}{|p{4.86cm}}{Chp 5 - Regulation} & 
\multicolumn{1}{|p{4.03cm}}{Chp 6 - Energetics} & 
\multicolumn{1}{|p{5.71cm}|}{Muscle Function Overview} \\ 
\hline
\multicolumn{1}{|p{1.9cm}}{June 23} & 
\multicolumn{1}{|p{4.86cm}}{Quiz 2: Chps 3-6} & 
\multicolumn{1}{|p{4.03cm}}{Chp 7 - Microcirculation} & 
\multicolumn{1}{|p{5.71cm}|}{Anatomy Lecture (Dr. Duffy)} \\ 
\hline
\multicolumn{1}{|p{1.9cm}}{June 30} & 
\multicolumn{1}{|p{4.86cm}}{Chp 8 – Renal Clearance} & 
\multicolumn{1}{|p{4.03cm}}{Chp 8 – Renal (cont)} & 
\multicolumn{1}{|p{5.71cm}|}{No Class – July 4\textsuperscript{th}} \\ 
\hline
\multicolumn{1}{|p{1.9cm}}{July 7} & 
\multicolumn{1}{|p{4.86cm}}{Quiz 3 – Chps 7 \& 8} & 
\multicolumn{1}{|p{4.03cm}}{Chp 9 – Circulation} & 
\multicolumn{1}{|p{5.71cm}|}{Chp 9 (cont) \& Chp 10 Cardiac Pump} \\ 
\hline
\multicolumn{1}{|p{1.9cm}}{July 14} & 
\multicolumn{1}{|p{4.86cm}}{Chp 10 – Cardiac Pump } & 
\multicolumn{1}{|p{4.03cm}}{Quiz 4 – Chps 9 \& 10 } & 
\multicolumn{1}{|p{5.71cm}|}{Chp 11 - Respiration} \\ 
\hline
\multicolumn{1}{|p{1.9cm}}{July 21} & 
\multicolumn{1}{|p{4.86cm}}{Chp 11 - Respiration (cont)} & 
\multicolumn{1}{|p{4.03cm}}{Chp 12 - Ventilation} & 
\multicolumn{1}{|p{5.71cm}|}{Chp 12 - Ventilation (cont)} \\ 
\hline
\multicolumn{1}{|p{1.9cm}}{July 28} & 
\multicolumn{1}{|p{4.86cm}}{Quiz 5 – Chps 11 \& 12} & 
\multicolumn{1}{|p{4.03cm}}{Vital Signs Lab} & 
\multicolumn{1}{|p{5.71cm}|}{Vital Signs Lab} \\ 
\hline\hline
\multicolumn{1}{|p{1.9cm}}{August 4} & 
\multicolumn{3}{|p{14.6cm}|}{\centering
Final Exam \& OSCE as Scheduled} \\ 
\hline\hline
\end{tabular}
\end{table}

\bigskip
\section{Mea Culpa, Veniam Peto}

\textit{Mea Culpa, Veniam Peto} is Latin for my fault, and I ask for forgiveness. It is an acknowledgment of a personal error or fault and a request to accept my apologies for any errors in this syllabus. They are bound to occur. Please bring any errors you find to my attention as soon as possible.

\end{document}
